\section{Complete binary counting}
\label{sec:binary-complete}
Let $F_N$ be the number of all fixed-point-free $2$-subgroups of
$S_{2N}$ divided by $L_{2N}=(2N)!c_NG_N$, and let $E_N$ be the
same normalized count restricted to groups with a noncritical orbit.
Thus $F_N=C_{2N}^{\crit}+E_N$ exactly. We prove the complete recurrence
needed in Section~\ref{sec:assembly}.

\subsection{The original degree-thirty-two boundary}
\begin{theorem}[Adaptive cover at degree sixteen]\label{thm:binary16-cover}
Let $T\leq S_{16}$ be transitive and binary, and let $A$ be a section
of its original permutation module whose action factors through a
specified quotient $B$ of $T$. There are $C\leq A^B$ and a specified
permutation cover of $B$ of degree $v_0\in\{8,16\}$ such that
\begin{equation}\label{eq:binary16-cover}
 v_0+2\dim C+4\dim(A/C)^B\leq30.
\end{equation}
\end{theorem}
\begin{proof}
Use $S,t,H,\mathcal L,\ell,H_0$ from
Lemma~\ref{lem:joint-budget}. Put $d_0=\dim H_0\leq4$.
For $a=4$ the joint budgets $2^jB_{4-j}$, for $0\leq j\leq4$,
are $(6,6,8,8,16)$. In particular $t,\ell\leq6$.
If $d_0\leq3$, taking $W=H$ gives $t+\ell\leq8$; taking a
character line gives $j_0\leq6-t$. The separator in
Lemma~\ref{lem:separator} supplies $q$ evaluations with kernel
$C\leq S$, and \eqref{eq:fixed-quotient} gives
$\dim C=t-q$, $\dim(A/C)^B=q$.
If $t\leq5$, their cost is
$2t+2q\leq t+\ell+6\leq14$.
If $t=6$, then $\ell\leq2$ and $j_0=0$, so $q\leq1$ and the
same bound holds. The original degree-sixteen cover $T\to B$
therefore proves \eqref{eq:binary16-cover}.

If $d_0=4$, the common kernel of $H_0$ has sixteen orbits on sixteen
points. Faithfulness makes it trivial. Thus $T=C_2^4$ is regular;
as every character in $H_0$ comes from $B$, also $B=T$.
Four independent characters give a faithful action of this same $B$
on four disjoint pairs. Use it as the degree-eight cover.
The original module still has $t,\ell\leq6$, $j_0\leq6-t$.
For $t\leq5$ the separator cost is at most $t+\ell+6\leq17$;
for $t=6$ it has $q\leq\lceil\ell/2\rceil\leq3$ and cost at most
$18$. Adding the new cover degree gives at most $26$.
The module $A$ was never replaced by a section of the eight-point
permutation module.
\end{proof}

For every transitive binary action $U$ of degree $32$ and every
$N\lhd U$, choose an invariant pair system with kernel $K$ and top
$T\leq S_{16}$. The groups
$Q=U/N$, $A=KN/N$ and $B=U/(KN)$ meet the theorem's hypotheses.
Since every simple binary module of a $2$-group is trivial,
$r(A/C)=\dim(A/C)^B$. Proposition~\ref{prop:prefix} gives the
same-source envelope with $v=v_0+2\dim C$, $\lambda=r(A/C)/2$ and
$g=32-v-8\lambda\geq2$.
The finite-menu proof of Theorem~\ref{thm:fusion}, restricted on its
physical side to fixed-point-free binary sources, consequently gives
\begin{equation}\label{eq:binary32-fusion}
 T_{32}(N)\leq\eta_{32}(N)+K_{32}(N,N-16)F_{N-16},\qquad
 \eta_{32}(N)=2^{-\Omega(N^2)},\quad
 K_{32}(N,N-16)=2^{-\Omega(N)}.
\end{equation}
Here $T_{32}$ denotes the family with a degree-$32$ orbit and none
larger. The auxiliary graph may have fixed points, so its moment is
bounded by the complete count $s_{b+qv}$ as in
Proposition~\ref{prop:prefix}. Only the physical cold complement is
fixed-point-free. At this fixed width all action, normal, frame and
cut choices form a finite menu, and may all be retained at a fixed
constant cost. The coefficient is the sum of
$\Delta(2N,2N-32)D_C2^{(4-g/16)(2N-32)}/a(U)$;
its logarithm is at most $-g(2N-32)/16+O(\log(N+2))$.

\subsection{Uniformity over unbounded binary widths}
\begin{theorem}[Complete large-orbit deletion]\label{thm:binary-unbounded}
Let $T_{64}(N)$ be the normalized family with some actual orbit of
width at least $64$. There are nonnegative $\eta_{64}$ and $K_{64}$,
with $K_{64}(N,M)=0$ unless $M\leq N-32$, such that eventually
\begin{equation}\label{eq:binary-unbounded}
 T_{64}(N)\leq\eta_{64}(N)+\sum_{M<N}K_{64}(N,M)F_M,
 \qquad \eta_{64}(N)\leq2^{-(2N)^2/2^{22}},\quad
 \sum_MK_{64}(N,M)\leq2^{-(2N)/512}.
\end{equation}
\end{theorem}
\begin{proof}
Every nontrivial transitive $2$-group has an invariant pair system:
a central involution is fixed-point-free, since a central element
fixing a point fixes its entire transitive orbit. An invariant pair
system is determined by the partner of one point, so there are at
most $w-1$ systems. Retain all of them.
Write $w=2s\geq64$. For every literal normal $N\lhd U$ and every
pair kernel $K$, put
\[
 Q=U/N,\qquad A=KN/N,\qquad B=U/(KN),\qquad S=A^B.
\]
For $s=32$ use Theorem~\ref{thm:binary32-capacity}, and for
$s\geq64$ use Theorem~\ref{thm:binary-large-capacity}.
Retain all cuts $C\leq S$ satisfying
$2c+4r(A/C)\leq15s/16$, where $c=\dim C$.
This menu is nonempty and invariant under the original action
normalizer. The original pair top is a specified degree-$s$ cover of
$B$. Proposition~\ref{prop:prefix} gives
\[
 e_Q(J)\leq D_C2^{\lambda b}\Phi_C(J),\qquad
 \Phi_C(J)=e_B(J)2^{c d_2(J)},\quad
 D_C=|A/C||H^1(B,A/C)|,\quad\lambda=r(A/C)/2.
\]
The module bound in Lemma~\ref{lem:joint-budget} gives
$t=\dim S\leq B_{\log s}\leq5s/16$. Therefore
\begin{equation}\label{eq:binary-wide-parameters}
 w/2\leq v=s+2c\leq13w/16,\qquad
 g=w-v-8\lambda\geq w/32,\qquad g\leq w-v\leq v.
\end{equation}

We next bound the entire varying menu. The soluble normal-generation
bound of \cite[Lemma 6.1]{RoneyDougalTracey2025} states
$d_U(N)\leq\zeta w/\sqrt{\log w}$ for every normal subgroup of a
soluble transitive $U\leq S_w$. Here all $U$ are $2$-groups.
Since $|U|\leq2^{w-1}$, padded normal-generating tuples bound their
number of normal subgroups by $2^{O(w^2/\sqrt{\log w})}$.
Applying the same theorem to $N=U$ bounds $d(U)$: normal generation
and ordinary generation agree for a finite $2$-group by its Frattini
quotient. Every such action is conjugate into a fixed Sylow subgroup
of $S_w$, and padded generating tuples there bound the number of
transitive action classes by the same expression.
Furthermore
\[
 t\leq B_{\log s}=O(s/\sqrt{\log s}),\qquad
 \log G_t=O(s^2/\log s),
\]
\[
 \log D_C\leq\dim(A/C)(1+d(B))=O(w^2/\sqrt{\log w}).
\]
The last inequality uses $d(B)\leq d(T)$ for the original transitive
top $T$. Together with the frame count and $1/a(U)\leq1$, these give
one absolute function $h(w)=O(w^2/\sqrt{\log w})$ with
\begin{equation}\label{eq:binary-wide-menu}
 \sum_{[U],\ \deg U=w}\frac1{a(U)}
        \sum_{\text{frames},N,C}D_C\leq2^{h(w)}.
\end{equation}
Every normal and every cut is included before taking this bound.

The graph proof of Proposition~\ref{prop:prefix} uses the same source
$J$ for all maps. For binary $J$ and binary covers its graph is binary,
but may have fixed points; in particular
$\sum_{J\leq S_b,\ J\text{ binary}}\Phi_C(J)^q\leq s_{b+qv}$.
Put $e=g/16$, $a_0=v/8+e$, and declare the entry hot if
$\Phi_C(J)>2^{a_0b}$. Use
$q=\max(1,\lceil8eb/v^2\rceil)$ in \eqref{eq:hot-square}.
For $b>0$ its rounding square is at most $v^2/16$, while
\eqref{eq:binary-wide-parameters} gives
\[
 g/v\geq1/26,\qquad b+qv\leq3b/2+v\leq3n/2.
\]
Consequently the main exponent of each hot term is at most
\begin{equation}\label{eq:wide-hot-exponent}
 -2^{-18}b^2-87w^2/4096.
\end{equation}
For $b=0$ the strict hot event is empty. The normalized hot sum has
the original factor
\[
 \frac{D_C}{b!a(U)c_{n/2}G_{n/2}}
       2^{\lambda b-(q-1)a_0b}s_{b+qv}.
\]
Thus the coarse subgroup bound contributes $O(n^{3/2})$ and
$c_{n/2}^{-1}$ contributes $O(n\log(n+2))$ to its logarithm;
dropping $1/b!\leq1$ is harmless for this upper bound.
For every $\delta>0$ there is $H_\delta<\infty$ such that
$h(w)\leq\delta w^2+H_\delta$ for all $w\geq64$.
Take $\delta=1/128$ in \eqref{eq:binary-wide-menu}.
The retained negative terms in \eqref{eq:wide-hot-exponent}, together
with $b^2+w^2\geq n^2/2$, absorb all errors and the at most $n$ widths,
yielding the stated hot scalar.

On the complement of the union of all hot events, the complete cold
coefficient for $b=n-w$ is
\[
 \Delta(n,b)\sum_{[U],\deg U=w}\frac1{a(U)}
       \sum_{\text{frames},N,C}D_C2^{(w/8-g/16)b}.
\]
The source factorial cancels the $b!$ in the pointing factor once.
The coefficient inequality $c_{b/2}/c_{n/2}\leq(2N)^{w/2}$ and
$\log(G_{b/2}/G_{n/2})=-wb/8-w^2/16+O(1)$ give its logarithm at most
\[
 h(w)+(w/2)\log n+O(1)-wb/512-w^2/16.
\]
Use $h(w)\leq w^2/128+H$ and, uniformly for $64\leq w\leq n$,
absorb $(w/2)\log n+H+O(1)$ into $wn/8192$ for large $n$.
The result is at most $-wn/1024$. Summing at most $n$ widths proves
\eqref{eq:binary-unbounded}, conservatively with rate $n/512$.

The hot union is selected first over all actual entries; its complement
is then bounded by cold pointing. Deleting an actual orbit leaves a
complete fixed-point-free binary subgroup, so the cold target is
precisely $F_{b/2}$. This establishes the inequality before any bound
on $F$ has been assumed.
\end{proof}

\subsection{Finite witness coverage and its universal interpretation}
Let $\mathcal A$ be the seventeen-action alphabet in
Theorem~\ref{thm:binary-mixture}.

\begin{lemma}[Finite coverage principles]\label{lem:binary-coverage}
Fix a power-of-two degree $w$. A list of literal transitive $2$-groups
containing a Sylow subgroup of $S_w$, and closed up to supplied
permutation conjugacy under all transitive index-two subgroups,
contains every transitive binary action up to conjugacy.
For any finite $2$-group $U$, a list of literal normal subgroups
containing $1$ and closed under all inverse images of order-two
subgroups of $Z(U/N)$ contains every normal subgroup of $U$.
\end{lemma}
\begin{proof}
The iterated binary wreath group on $w$ points has order $2^{w-1}$,
which is the full $2$-part of $w!$. Hence every binary action is
conjugate into this explicit Sylow group. In a finite $2$-group any
subgroup lies at the bottom of a chain of index-two subgroups: choose
successive maximal proper overgroups. Every overgroup of a transitive
subgroup is transitive. Induction down the chain proves the first claim.
All index-two subgroups are exactly the kernels of the nonzero
functionals on $U/\Phi(U)$, so this closure can be checked exhaustively.

If $N<L$ are normal in $U$, the nontrivial normal subgroup $L/N$ of
$U/N$ meets its centre nontrivially, by the class equation. This
intersection contains an involution. Its inverse image is a normal
index-two extension of $N$ lying inside $L$. Repetition starting at
$1$ reaches every $L$, proving the second claim.
\end{proof}

\begin{lemma}[Finite binary entry theorem]\label{lem:binary-finite-menu}
For every transitive binary action $U$ of width $w\in\{2,4,8,16\}$
outside $\mathcal A$, and every literal normal $N\lhd U$, at least
one of the following witnesses exists for $Q=U/N$:
\begin{enumerate}
\item An invariant pair system with kernel $K$, an actual central cut
$C\leq(KN/N)^{U/(KN)}$, and an actual permutation cover of $B=U/(KN)$,
such that $v_B+2\dim C+4\dim((KN/N)/C)^B<w$.
\item A faithful permutation representation of $Q$ of degree $v<w$.
\item With $z=\dim\Omega_1Z(Q)$, the exact integer inequality
$38^{8z}<2^w25^{8z}$.
\item An equal-degree full-subdirect replacement by a word in
$\mathcal A$, with actual maps onto $Q$, retaining at least one
quarter of its points in noncritical factors.
\end{enumerate}
The last branch needs only the following four entries:
\[
\begin{array}{c|c|c|l}
 U&|N|&Q&\text{replacement word}\\\hline
 8T16&2&\operatorname{SmallGroup}(16,3)&8T27\\
 8T20&2&\operatorname{SmallGroup}(16,3)&8T27\\
 8T21&2&\operatorname{SmallGroup}(16,3)&8T27\\
 16T1086&4&\operatorname{SmallGroup}(256,7667)&C_4[4],D_8[4]^3
\end{array}
\]
In the last row the cover is a proper full-subdirect subgroup of order
$256$ in the displayed product, not that whole product.
\end{lemma}
\begin{proof}
This is the finite computer-assisted lemma recorded by the literal
binary entry certificate. We describe its mathematical coverage and
the statements checked at each vertex; these are essential to its use.
There are explicit iterated wreath roots in each of the four degrees.
Each action record gives every transitive index-two child, an actual
conjugator to its destination, and a checked basis of its Frattini
quotient. Nonregular base actions retain their children. Regular
transitive actions have no transitive proper subgroup.
Each nonbase action record gives $1$ and every central-involution
extension of each normal vertex. The two closure checks prove exhaustive
action and normal coverage by Lemma~\ref{lem:binary-coverage}; catalogue
identifiers play no role in these implications.

For the first branch the record supplies the pair partition, kernel
$K$, intersection $M=K\cap N$ and lifted cut $M\leq C_0\leq K$.
The literal equalities $[C_0,U]\leq M$ and
$\ker(U\to\text{cover image})=KN$ certify the same central section
and quotient. Fixed sections are computed by commutators modulo
$M$ and $C_0$, giving the integers $c,r$ and the strict gap.
For a direct witness the supplied generator images define a map with
kernel exactly $N$ and moved degree strictly below $w$.
For a character witness the centre and its involution subgroup are
reconstructed in the literal quotient, and the displayed integer
inequality is checked. For a transport witness both specified maps
are onto the same literal quotient, the original map has kernel $N$,
and each displayed projection of the actual cover group equals its
assigned action. The degree and noncritical fraction are then checked
on the actual orbit partition. The four exceptional charts satisfy
these conditions, including the retained proper cover in degree sixteen.
The committed finite records and checker in
Appendix~\ref{sec:certificates} supply these finite identities.

There is also a uniform proof for regular actions of width $w\geq8$:
choose a central involution and its coset pairs. Then $K=C_2$,
$A=KN/N$ is central of dimension at most one, and one may take $C=A$.
The quotient $B$ has its regular cover of degree at most $w/2$;
$w/2+2\dim A<w$. All binary actions of degrees two and four are
already in $\mathcal A$.
\end{proof}

The first two witness types have the graph envelopes of
Proposition~\ref{prop:prefix}, with $g>0$. The third has a direct
slope strictly below $w/8$. Indeed a $2$-group $Q$ has a faithful
tuple of $z$ irreducible complex characters: choose constituents whose
central characters span the dual of $\Omega_1Z(Q)$. A nontrivial
normal subgroup meets $\Omega_1Z(Q)$, so their common kernel is trivial.
For a binary $J\leq S_b$, inflation along an epimorphism yields a
tuple from $\operatorname{Irr}(J)^z$. Every tuple determines at most
one kernel, and each kernel has exactly $|\Aut Q|$ epimorphisms to $Q$.
The nilpotent conjugacy-class bound
\cite[Theorem 1.5]{Maroti2005Classes} therefore gives
\begin{equation}\label{eq:binary-character-entry}
 |\Epi(J,Q)|\leq|\Aut Q|\,k(J)^z
 \leq|\Aut Q|\,2^{z\log(38/25)b}.
\end{equation}
The strict integer test is exactly $z\log(38/25)<w/8$.
For $Q=1$ use $e_Q(J)=1$. These assertions concern the actual source
$J$; no faithful action of an abstract quotient on its points is assumed.

\subsection{Small widths and the complete recurrence}
\begin{theorem}[Complete binary recurrence]\label{thm:binary-recurrence}
There are nonnegative $\eta_B(N)$ and $K_B(N,M)$, supported on
$0\leq M<N$, and constants $a_B,C_B>0$, such that eventually
\begin{equation}\label{eq:proved-binary-recurrence}
 E_N\leq\eta_B(N)+\sum_{M<N}K_B(N,M)F_M,\qquad
 \eta_B(N)+\sum_{M<N}K_B(N,M)\leq C_B2^{-a_BN}.
\end{equation}
In particular $F_N$ is bounded and $E_N=O(2^{-a_BN})$.
\end{theorem}
\begin{proof}
Power-of-two orbit sizes give the exact disjoint partition
\[
 E_N=T_{64}(N)+T_{32}(N)+S_{16}(N),
\]
where $S_{16}$ has all actual orbit widths at most sixteen and at least
one noncritical orbit. The first two parts are bounded by
\eqref{eq:binary-unbounded} and \eqref{eq:binary32-fusion}.

For $S_{16}$ consider a nonbase orbit $U$ and its actual axis
$N=H\cap U$. Include a fixed finite invariant menu of all witness
transports in Lemma~\ref{lem:binary-finite-menu}; choose the transports
before counting the complementary source. First take the union in
which some actual graph entry is hot. Theorem~\ref{thm:fusion}
bounds this complete union by a scalar $2^{-\Omega(N^2)}$.
On its complement, point any orbit with an accepted strict local
entry. For graph entries the cold coefficient is
\[
 \Delta(2N,2N-w)\frac{D}{a(U)}
      2^{(w/8-g/16)(2N-w)}.
\]
For the direct character entry it is the same coefficient with
$D=|\Aut Q|$ and exponent $z\log(38/25)(2N-w)$.
Their original Gaussian ratio leaves respectively
$-g(2N-w)/16$ or
$-(w/8-z\log(38/25))(2N-w)$, with error $O(\log(N+2))$.
The fixed finite menu has a positive minimum gap. Summing every literal
normal and every witness thus yields a kernel $K_s$ with row
$2^{-\Omega(N)}$, supported at $N-4,N-8$; only nonnegative targets
are included. Entries with zero prefix degree have a direct bound and
need no hot moment. Deleting the original orbit again leaves the
complete fixed-point-free binary count $F_{N-w/2}$.

In the remaining family every nonbase orbit has a transport witness
for its exact axis. Apply all these maps simultaneously, leaving
$\mathcal A$ factors untouched. The product of the embedded axes lies
in $H$, so Proposition~\ref{prop:transport} recovers $H$ as the full
preimage of the joint quotient image. Every replacement has the same
physical degree, and retains at least one quarter of its noncritical
support. Theorem~\ref{thm:binary-mixture} and the stronger profile
envelopes \eqref{eq:proved-mixture-envelopes}, inserted in
Proposition~\ref{prop:absorb-decorations}, give a scalar
$2^{-\Omega(N)}$ for the entire remaining family. This also includes
words already contained in $\mathcal A$, by identity transport.
The small original noncritical-support range is removed before
charging the decoration factor.

Consequently
$S_{16}(N)\leq\eta_s(N)+\sum_MK_s(N,M)F_M$, with both its scalar
and complete row exponentially small. Add this to the two disjoint
large-width bounds to obtain \eqref{eq:proved-binary-recurrence}.
No estimate of $F_M$ has been used to prove the kernel rows.

Finally $C_{2M}^{\crit}$ is bounded by the independently proved
critical-family theorem. Substitute
$F_M=C_{2M}^{\crit}+E_M$ in the recurrence. Beyond an onset its row
is at most $1/2$, and its forcing is bounded. A constant dominating
twice that forcing and all finitely many initial $E_M$ bounds $E_N$
by induction. Thus $F_N$ is bounded. Substituting this established
bound back into the same recurrence proves exponential decay.
\end{proof}
